% Part: incompleteness
% Chapter: representability-in-q
% Section: introduction

\documentclass[../../../include/open-logic-section]{subfiles}

\begin{document}

\olfileid{inc}{req}{int}
\olsection{పరిచయం}

గణనీయ ప్రమేయాల గురించిన ప్రాథమిక వాస్తవాలను వ్యక్తం చేసి నిరూపించగల
సిద్ధాంతాలకు అసంపూర్ణతా సిద్ధాంతాలు వర్తిస్తాయి. అలాంటి అత్యంత మితమైన ఒక
సిద్ధాంతాన్ని ``$\Th{Q}$'' (లేదా రాఫెల్ రాబిన్సన్ పేరు మీద కొన్నిసార్లు
``రాబిన్సన్~$Q$'') అని పిలిచి వివరిస్తాం. ఒక ప్రమేయం~$\Th{Q}$లో
\emph{ప్రాతినిధ్యం చేయదగినది} అనడం అంటే ఏమిటో చెప్పి, తరువాత కింది ఫలితాన్ని
నిరూపిస్తాం:
\begin{quote}
  ఒక ప్రమేయం~$\Th{Q}$లో ప్రాతినిధ్యం చేయదగినదైతే, అప్పుడు మాత్రమే అది
  గణనీయమైనది.
\end{quote}
ఇది ఒక వైపు గణనీయతకు మరొక నమూనాను ఇస్తుంది. మరోవైపు, ఆగే సమస్యను దానికి
తగ్గించడం ద్వారా $\Setabs{!A}{\Th{Q} \Proves !A}$ సమితి నిర్ణయించదగినది
కాదని చూపడానికి కూడా దీన్ని వాడతాం. ఈ చర్చ ముగిసేసరికి దీనికంటే చాలా
బలమైన ఫలితాలను నిరూపించి ఉంటాం.

$\Th{Q}$ భాష అంకగణిత భాష; $\Th{Q}$ కింది స్వీకృతాలను కలిగి ఉంటుంది
(వీటిని \tetoken{తాదాత్మ్య విధేయంతో}{!!{identity}} కూడిన మొదటిస్థాయి తర్కపు ఇతర స్వీకృతాలు, నియమాలతో
కలిపి వాడాలి):
\begin{align*}
& \lforall[x][\lforall[y][(\eq[x'][y'] \lif \eq[x][y])]] \tag{$!Q_1$}\\
& \lforall[x][\eq/[\Obj 0][x']] \tag{$!Q_2$}\\
& \lforall[x][(\eq[x][\Obj 0] \lor \lexists[y][\eq[x][y']])] \tag{$!Q_3$}\\
& \lforall[x][\eq[(x + \Obj 0)][x]] \tag{$!Q_4$}\\
& \lforall[x][\lforall[y][\eq[(x + y')][(x + y)']]] \tag{$!Q_5$}\\
& \lforall[x][\eq[(x \times \Obj 0)][\Obj 0]] \tag{$!Q_6$}\\
& \lforall[x][\lforall[y][\eq[(x \times y')][((x \times y) + x)]]] \tag{$!Q_7$}\\
& \lforall[x][\lforall[y][(x < y \liff \lexists[z][\eq[(z' + x)][y]])]] \tag{$!Q_8$}
\end{align*}
ప్రతి సహజ సంఖ్య~$n$కు, మొత్తం $n$ టిక్ గుర్తులు ఉండే
$\Obj{0}^{\prime\prime\ldots\prime}$ పదాన్ని సంఖ్యాపదం~$\num{n}$గా
నిర్వచించండి. అందువల్ల $\num{0}$ అనేది \tetoken{స్థిరసంకేతం}{!!{constant}}~$\Obj{0}$ ఒక్కటే;
$\num{1}$ అనేది $\Obj{0}'$, $\num{2}$ అనేది $\Obj{0}''$, ఇలాగే కొనసాగుతుంది.

అంకగణిత సిద్ధాంతంగా $\Th{Q}$ \emph{అత్యంత} బలహీనమైనది. ఉదాహరణకు,
$\lforall[x][\eq/[x][x']]$ లేదా $\lforall[x][\lforall[y][(x + y) = (y +
    x)]]$ వంటి అతి సరళమైన వాస్తవాలను కూడా అందులో నిరూపించలేం. అయితే
$\Th{Q}$ అంత ఆసక్తికరంగా ఉండటానికి అది ఇంత బలహీనంగా ఉండడమే చాలావరకు కారణమని
చూస్తాం. నిజానికి అసంపూర్ణతా సిద్ధాంతం వర్తించడానికి అవసరమైన కనిష్ఠ బలానికి
అది అతి సమీపంలో ఉంది. $\Th{Q}$కు \emph{పరిమిత} స్వీకృత సమితి ఉండడం కూడా
దాని ఆసక్తికి మరొక కారణం.

$\Th{Q}$కు ఆగమన పథకాన్ని చేరిస్తే, $\Th{Q}$కంటే బలమైన \emph{పియానో
అంకగణితం}~$\Th{PA}$ అనే సిద్ధాంతం వస్తుంది:
\[
(!A(\Obj 0) \land \lforall[x][(!A(x) \lif !A(x'))]) \lif \lforall[x][!A(x)]
\]
ఇక్కడ $!A(x)$ ఏ సూత్రమైనా కావచ్చు. $!A(x)$లో $x$ కాకుండా ఇతర స్వేచ్ఛా
\tetoken{చరాలు}{!!{variable}s} ఉంటే, వాటన్నింటినీ బంధించడానికి ముందు సార్వత్రిక పరిమాణకాలను
చేర్చుతాం (దానివల్ల ఆగమన పథకపు అనుగుణ రూపం \tetoken{ఒక వాక్యం}{!!a{sentence}} అవుతుంది).
ఉదాహరణకు, $!A(x,y)$లో \tetoken{చరం}{!!{variable}}~$y$ కూడా స్వేచ్ఛగా ఉంటే, అనుగుణ రూపం
\[
\lforall[y][((!A(\Obj 0) \land \lforall[x][(!A(x) \lif !A(x'))]) \lif
  \lforall[x][!A(x)])]
\]
$\Th{Q}$ స్వీకృతాల నుంచి నిరూపించగలిగినదానికంటే ఆగమన పథక రూపాలను వాడి
$\Th{PA}$ స్వీకృతాల నుంచి ఎంతో ఎక్కువను నిరూపించవచ్చు. నిజానికి పియానో
అంకగణితంలో నిరూపించలేని సహజ సంఖ్యల గురించిన ``సహజమైన'' వాక్యాలను కనుగొనడానికే
చాలా శ్రమ అవసరం!{}

\begin{defn}
\ollabel{defn:representable-fn}
  సహజ సంఖ్యల నుంచి సహజ సంఖ్యలకు వెళ్లే ప్రమేయం
  $f(x_0,\ldots,x_k)$ను కింది షరతు నెరవేరితే~$\Th{Q}$లో
  \emph{ప్రాతినిధ్యం చేయదగినది} అంటాం: ఏ
  $f(n_0,\dots,n_k)=m$కైనా $\Th{Q}$ కింది రెండింటినీ నిరూపించేలా ఒక
  సూత్రం~$!A_f(x_0,\dots,x_k,y)$ ఉండాలి:
\begin{enumerate}
\item\ollabel{defn:rep:a} $!A_f(\num{n_0}, \dots, \num{n_k}, \num{m})$
\item\ollabel{defn:rep:b} $\lforall[y][(!A_f(\num{n_0}, \dots,
\num{n_k}, y) \lif \num{m} = y)]$.
\end{enumerate}
\end{defn}

ఈ నిర్వచనాన్ని చెప్పడానికి ఇతర మార్గాలూ ఉన్నాయి. ఉదాహరణకు, $\Th{Q}$
$\lforall[y][(!A_f(\num{n_0}, \dots, \num{n_k}, y) \liff
\eq[y][\num{m}])]$ను నిరూపించాలని సమానార్థకంగా కోరవచ్చు.

\begin{thm}
\ollabel{thm:representable-iff-comp}
ఒక ప్రమేయం~$\Th{Q}$లో ప్రాతినిధ్యం చేయదగినదైతే, అప్పుడు మాత్రమే అది
గణనీయమైనది.
\end{thm}

ఈ సిద్ధాంతాన్ని నిరూపించడంలో రెండు దిశలున్నాయి. వాక్యనిర్మాణపు
అంకగణితీకరణ సిద్ధమైన తరువాత ఎడమ-నుంచి-కుడికి దిశ చాలావరకు సూటిగా ఉంటుంది.
మరో దిశకు ఎక్కువ పని అవసరం. ప్రాథమిక ఆలోచన ఇది: ``గణనీయమైనది'' అనే భావాన్ని
ఖచ్చితం చేయడానికి ``సామాన్య పునరావృత్తం''ను ఎంచుకుని, ప్రతి సామాన్య పునరావృత్త
ప్రమేయం~$\Th{Q}$లో ప్రాతినిధ్యం చేయదగినదని చూపుతాం. ఒక ప్రమేయాన్ని~$\Zero$,
ఉత్తరగామి ప్రమేయం~$\Succ$, ప్రక్షేప ప్రమేయాలు~$\Proj{n}{i}$ నుంచి సంయుక్తం,
ఆదిమ పునరావృత్తి, సక్రమ కనిష్ఠీకరణలను వాడి నిర్వచించగలిగితే, అది సామాన్య
పునరావృత్త ప్రమేయమని గుర్తుంచుకోండి. కాబట్టి ప్రతి సామాన్య పునరావృత్త
ప్రమేయం~$\Th{Q}$లో ప్రాతినిధ్యం చేయదగినదని చూపడానికి ఒక మార్గం: ప్రాథమిక
ప్రమేయాలు ప్రాతినిధ్యం చేయదగినవని, అలాగే ప్రాతినిధ్యం చేయదగిన ప్రమేయాల నుంచి
సంయుక్తం, ఆదిమ పునరావృత్తి, సక్రమ కనిష్ఠీకరణల ద్వారా నిర్వచించిన ప్రమేయాలు
కూడా ప్రాతినిధ్యం చేయదగినవని చూపడం. మరో మాటలో, ప్రాథమిక ప్రమేయాలు
ప్రాతినిధ్యం చేయదగినవనీ, ప్రాతినిధ్యం చేయదగిన ప్రమేయాల తరగతి సంయుక్తం,
ఆదిమ పునరావృత్తి, సక్రమ కనిష్ఠీకరణల కింద ``సంవృతం'' అనీ చూపవచ్చు. ఇది ప్రతి
సామాన్య పునరావృత్త ప్రమేయమూ ప్రాతినిధ్యం చేయదగినదని హామీ ఇస్తుంది.

అయితే ప్రాతినిధ్యం చేయదగిన ప్రమేయాలు ఆదిమ పునరావృత్తి కింద సంవృతమని
చూపాల్సిన మెట్టు కష్టమైనదని తేలుతుంది. ఆ మెట్టును తప్పించుకోవడానికి మొదట
ఆదిమ పునరావృత్తి లేకుండానే పని చేయవచ్చని చూపుతాం. అంటే ప్రతి సామాన్య
పునరావృత్త ప్రమేయాన్ని ప్రాథమిక ప్రమేయాల నుంచి సంయుక్తం, సక్రమ కనిష్ఠీకరణ
మాత్రమే వాడి నిర్వచించవచ్చని చూపుతాం. ఇందుకోసం ఆదిమ పునరావృత్తిని ఒక
నిర్దిష్ట సక్రమ కనిష్ఠీకరణ ద్వారానే చేయవచ్చని చూపుతాం. అయితే ఇది పనిచేయాలంటే
కూడిక, గుణకారం, సర్వసమానతా సంబంధపు లక్షణ ప్రమేయం~$\Char{=}$ అనే కొన్ని అదనపు
ప్రాథమిక ప్రమేయాలను చేర్చాలి. తరువాత ఈ ప్రాథమిక ప్రమేయాలన్నీ~$\Th{Q}$లో
ప్రాతినిధ్యం చేయదగినవని, ప్రాతినిధ్యం చేయదగిన ప్రమేయాలు సంయుక్తం, సక్రమ
కనిష్ఠీకరణల కింద సంవృతమని చూపడం ద్వారా సిద్ధాంతాన్ని నిరూపించవచ్చు.

\end{document}
